\chapter{附录：待补实验与写作占位}

\section{Layer B 必补信息}

\begin{itemize}[leftmargin=2em]
  \item 板卡型号、SoC 或 stepping
  \item Linux 内核版本与 \code{uname -a}
  \item \capstone\ 版本与提交哈希
  \item 实验命令、运行时间与日志目录
  \item 至少三次复跑结果
\end{itemize}

\section{Layer A / Layer B 对照建议}

\begin{itemize}[leftmargin=2em]
  \item 固定扫描范围或固定 seed
  \item 明确启用的扩展集合与过滤选项
  \item 将差异分桶为模拟器、解码器、特权或状态依赖差异
  \item 对代表性样例保留 second-decoder 与人工 triage 记录
\end{itemize}

\section{图表回填清单}

当前 \code{paper/} 目录中仍保留以下正式占位。
\begin{itemize}[leftmargin=2em]
  \item 图 1-1：总体框架图
  \item 图 4-1：\memcage\ 内存布局与信号路径
  \item 图 4-2：\ptracebackend\ 工作流
  \item 图 6-1：Layer A/B/C 证据映射图
  \item 表 2-1：相关工作对比表
  \item 表 6-4：Layer B 实机复现实验表
  \item 表 6-5：Layer A / Layer B 差异分桶表
\end{itemize}
